% Reviewed translation: PDF pages 356--360 / printed pages 342--346.
\MathBookSourcePage{356}
\setcounter{statement}{0}
\begin{Lemma}\label{lem:navier-upwind-projection-mesh-norm}
假设 $\Omega$ 是有界凸多边形, 且 $\mathcal T_h$ 是 $\bar\Omega$ 的一致正则
三角剖分族. 设某个整数 $k$ 满足 $1\leq k\leq l$, 并且
$\phi\in H^{k+1}(\Omega)\cap H_0^1(\Omega)$. 则有估计
\begin{equation}\label{eq:navier-upwind-projection-mesh-estimate}
[\phi-\mathring P_h\phi]_t
\leq C_t h^{k-1}\|\phi\|_{k+1,\Omega}
\quad\forall t\geq2,
\end{equation}
其中常数 $C_t$ 与 $h$ 无关.
\end{Lemma}

这个 Lemma 是 Theorem~\ref{thm:appendix-ritz-projection-lp-error} 和下式的直接推论:
\begin{equation}\label{eq:navier-upwind-inverse-mesh-norm}
[\phi_h]_t\leq C(t)h^{-1}|\phi_h|_{1,\Omega}
\quad\forall\phi_h\in\Phi_h.
\end{equation}

这些准备工作可以导出
$\widetilde a(\mathord\cdot;\mathord\cdot,\mathord\cdot)$ 的若干基本性质.
首先, $\widetilde a(\mathord\cdot;\mathord\cdot,\mathord\cdot)$ 是有界的.

\setcounter{statement}{1}
\begin{Lemma}\label{lem:navier-upwind-form-bounded}
保留 Lemma~\ref{lem:navier-upwind-projection-mesh-norm} 的假设. 对所有
$v_h\in V_h$ 以及所有
$z_h$, $u_h=(\boldsymbol u_h=\operatorname{curl}\psi_h,\omega_h)$,
$w_h=(\boldsymbol w_h=\operatorname{curl}\chi_h,\xi_h)\in X_h$, 形式
$\widetilde a(\mathord\cdot;\mathord\cdot,\mathord\cdot)$ 满足
\begin{equation}\label{eq:navier-upwind-form-bound}
|a_1(u_h;v_h,\boldsymbol w_h)|
+|a_2^{z_h}(u_h;v_h,\boldsymbol w_h)|
\leq C\|\boldsymbol u_h\|_{0,4,\Omega}|v_h|
\|\boldsymbol w_h\|_{0,4,\Omega}.
\end{equation}
\end{Lemma}
\begin{Proof}
设 $v_h=(\operatorname{curl}\phi_h,\theta_h)\in V_h$.
由 Hölder 不等式,
\[
|a_1(u_h;v_h,\boldsymbol w_h)|
\leq C_1\|\psi_h\|_{1,4,\Omega}\|\chi_h\|_{1,4,\Omega}
\left(\sum_{\kappa\in\mathcal T_h}|\phi_h|_{2,\kappa}^2\right)^{1/2}.
\]
因此 \eqref{eq:navier-upwind-uniform-broken-h2} 给出
\[
|a_1(u_h;v_h,\boldsymbol w_h)|
\leq C_2\|\boldsymbol u_h\|_{0,4,\Omega}|v_h|
\|\boldsymbol w_h\|_{0,4,\Omega}.
\]

接着考虑 $a_2^{z_h}(u_h;Pv_h,\boldsymbol w_h)$, 其中
$Pv_h=(\operatorname{curl}\phi,\theta_h)$. 由于 $\phi\in H^2(\Omega)$,
$a_2(\mathord\cdot;\mathord\cdot,\mathord\cdot)$ 中的面积分在
$\mathcal T_h$ 的所有内部边 $\kappa'$ 上都消失; 又由于 $u_h\in X_h$,
因子 $\operatorname{curl}\psi_h\cdot\boldsymbol n$ 在 $\mathcal T_h$ 的所有边界边上消失.
因此
\begin{equation}\label{eq:navier-upwind-smoothed-boundary-zero}
a_2^{z_h}(u_h;Pv_h,\boldsymbol w_h)=0,
\end{equation}
从而
\[
|a_2^{z_h}(u_h;v_h,\boldsymbol w_h)|
=|a_2^{z_h}(u_h;v_h-Pv_h,\boldsymbol w_h)|.
\]
于是
\[
\begin{aligned}
|a_2^{z_h}(u_h;v_h,\boldsymbol w_h)|
\leq C_3
&\left(\sum_{\kappa\in\mathcal T_h}
\|\operatorname{curl}\psi_h\|_{0,4,\partial\kappa}^4\right)^{1/4}\\
{}\times
&\left(\sum_{\kappa\in\mathcal T_h}
\|\operatorname{curl}\chi_h\|_{0,4,\partial\kappa}^4\right)^{1/4}
\left(\sum_{\kappa\in\mathcal T_h}
\|\operatorname{curl}(\phi_h-\phi)\|_{0,\partial\kappa}^2\right)^{1/2}.
\end{aligned}
\]
一方面, 例行应用
\[
\|\theta\|_{0,t,\partial\kappa}
\leq C_4(h_\kappa^{1/t}/\rho_\kappa)
(\|\theta\|_{0,\kappa}^2+h_\kappa^2|\theta|_{1,\kappa}^2)^{1/2}
\quad\forall\theta\in H^1(\kappa),\quad\forall t\geq1,
\]

\MathBookSourcePage{357}
给出
\[
\left(\sum_{\kappa\in\mathcal T_h}
\|\operatorname{curl}(\phi_h-\phi)\|_{0,\partial\kappa}^2\right)^{1/2}
\leq C_5h^{-1/2}\left(
|\phi_h-\phi|_{1,\Omega}
+h\left(\sum_{\kappa\in\mathcal T_h}|\phi_h-\phi|_{2,\kappa}^2\right)^{1/2}
\right).
\]
另一方面,
\[
\|\operatorname{curl}\theta_h\|_{0,4,\partial\kappa}
\leq C_6h^{-1/4}|\theta_h|_{1,4,\kappa}
\quad\forall\kappa\in\mathcal T_h,
\quad\forall\theta_h\in\Theta_h.
\]
把这些不等式与 \eqref{eq:navier-upwind-mesh-dependent-norm} 在 $t=2$ 时结合, 得到
\begin{equation}\label{eq:navier-upwind-boundary-form-estimate}
|a_2^{z_h}(u_h;v_h,\boldsymbol w_h)|
\leq C_7\|\boldsymbol u_h\|_{0,4,\Omega}
\|\boldsymbol w_h\|_{0,4,\Omega}[\phi_h-\phi]_2.
\end{equation}
\end{Proof}

下一个 Lemma 表明, 关于 $z_h$,
$a_2^{z_h}(\mathord\cdot;\mathord\cdot,\mathord\cdot)$ ``几乎 Lipschitz 连续''.

\setcounter{statement}{2}
\begin{Lemma}\label{lem:navier-upwind-boundary-almost-lipschitz}
保留 Lemma~\ref{lem:navier-upwind-form-bounded} 的记号和假设, 但假设
$v_h,w_h\in V_h$. 对 $X_h$ 中所有成对的
$z_h=(\boldsymbol z_h,\zeta_h)$ 和
$z_h^*=(\boldsymbol z_h^*,\zeta_h^*)$, 差
$a_2^{z_h}-a_2^{z_h^*}$ 满足
\begin{equation}\label{eq:navier-upwind-boundary-almost-lipschitz}
\begin{aligned}
&|a_2^{z_h}(u_h;v_h,\boldsymbol w_h)
-a_2^{z_h^*}(u_h;v_h,\boldsymbol w_h)|\\
&\qquad\leq Ch^{1/2}|v_h||w_h|
\begin{cases}
\|\boldsymbol u_h\|_{0,4,\Omega},&z_h^*\neq u_h,\\
\|\boldsymbol z_h-\boldsymbol z_h^*\|_{0,4,\Omega},&z_h^*=u_h,
\end{cases}
\end{aligned}
\end{equation}
对所有 $u_h,z_h,z_h^*\in X_h$ 以及所有 $v_h,w_h\in V_h$ 成立.
\end{Lemma}
\begin{Proof}
此证明基于恒等式
\begin{equation}\label{eq:navier-upwind-inflow-difference-identity}
\begin{aligned}
&a_2^{z_h}(u_h;v_h,\boldsymbol w_h)
-a_2^{z_h^*}(u_h;v_h,\boldsymbol w_h)\\
&\quad=\sum_{\kappa\in\mathcal T_h}
\int_{\partial_-\kappa(z_h-z_h^*)}
\boldsymbol u_h\cdot\boldsymbol n
(v_{hi}^{\mathrm{ext}}-v_{hi}^{\mathrm{int}})
(w_{hi}^{\mathrm{int}}-w_{hi}^{\mathrm{ext}})\,ds,
\end{aligned}
\end{equation}
其中
\[
\partial_-\kappa(z-z^*)
=\{x\in\partial\kappa;
\ \boldsymbol z\cdot\boldsymbol n(x)<0
\text{ and }\boldsymbol z^*\cdot\boldsymbol n(x)>0\}.
\]
因此, 根据 \eqref{eq:navier-upwind-smoothed-boundary-zero}, 不仅可把 $v_h$ 换成
$v_h-Pv_h$, 也可把 $w_h$ 换成 $w_h-Pw_h$. 于是
\[
\begin{aligned}
&|a_2^{z_h}(u_h;v_h,\boldsymbol w_h)
-a_2^{z_h^*}(u_h;v_h,\boldsymbol w_h)|\\
&\quad\leq C_1
\left(\sum_{\kappa\in\mathcal T_h}
\|\operatorname{curl}\psi_h\|_{0,4,\partial\kappa}^4\right)^{1/4}
\left(\sum_{\kappa\in\mathcal T_h}
\|\operatorname{curl}(\phi_h-\phi)\|_{0,4,\partial\kappa}^4\right)^{1/4}\\
&\qquad\times
\left(\sum_{\kappa\in\mathcal T_h}
\|\operatorname{curl}(\chi_h-\chi)\|_{0,\partial\kappa}^2\right)^{1/2}.
\end{aligned}
\]
前一 Lemma 的方法容易给出

\MathBookSourcePage{358}
\begin{equation}\label{eq:navier-upwind-inflow-difference-bound}
|a_2^{z_h}(u_h;v_h,\boldsymbol w_h)
-a_2^{z_h^*}(u_h;v_h,\boldsymbol w_h)|
\leq C_2h^{1/2}\|\boldsymbol u_h\|_{0,4,\Omega}
[\phi_h-\phi]_2[\chi_h-\chi]_2.
\end{equation}
因此, 一般情形下由 Lemma~\ref{lem:navier-upwind-projection-mesh-norm} 和
\eqref{eq:navier-upwind-smoothing-regularity} 得到
\eqref{eq:navier-upwind-boundary-almost-lipschitz}.

为处理 $z_h^*=u_h$ 的特殊情形, 注意
\[
|\boldsymbol u\cdot\boldsymbol n(x)|
\leq|\boldsymbol u\cdot\boldsymbol n(x)-\boldsymbol z\cdot\boldsymbol n(x)|
\quad\forall x\in\partial_-\kappa(z,-u).
\]
因此, \eqref{eq:navier-upwind-inflow-difference-identity} 右端中的因子
$\boldsymbol u_h\cdot\boldsymbol n$ 可由差
$|(\boldsymbol u_h-\boldsymbol z_h)\cdot\boldsymbol n|
=|(\boldsymbol z_h^*-\boldsymbol z_h)\cdot\boldsymbol n|$ 控制.
容易验证, 作此修改后上述论证的其余部分仍然有效.
\end{Proof}

最后, 还可证明
$\widetilde a(\mathord\cdot;\mathord\cdot,\mathord\cdot)$ 的另一个有用界.

\setcounter{statement}{3}
\begin{Lemma}\label{lem:navier-upwind-consistency-bound}
保留 Lemma~\ref{lem:navier-upwind-boundary-almost-lipschitz} 的假设和记号,
但取 $v_h\in X_h$. 对所有
$z=(\operatorname{curl}\alpha,\beta)$, 其中
$\alpha\in H^2(\Omega)\cap H_0^1(\Omega)$ 且 $\beta\in L^2(\Omega)$, 有估计
\begin{equation}\label{eq:navier-upwind-consistency-bound}
|\widetilde a(u_h;v_h-z,\boldsymbol w_h)|
\leq Ch^{1/2}\|\boldsymbol u_h\|_{0,4,\Omega}
[\phi_h-\alpha]_4|w_h|
\quad\forall u_h,v_h\in X_h,
\quad\forall w_h\in V_h,
\end{equation}
其中 $a_2(\mathord\cdot;\mathord\cdot,\mathord\cdot)$ 中的线积分可以取在
$\partial\kappa$ 的任意部分上.
\end{Lemma}
\begin{Proof}
取 $Pw_h=(\boldsymbol w=\operatorname{curl}\chi,\xi_h)$, 并令
$\boldsymbol z=\operatorname{curl}\alpha$. 利用 $\widetilde a$ 对最后一个自变量的线性性可写成
\[
\widetilde a(u_h;v_h-z,\boldsymbol w_h)
=\widetilde a(u_h;v_h-z,\boldsymbol w_h-\boldsymbol w)
+\widetilde a(u_h;v_h-z,\boldsymbol w).
\]
由于 $\boldsymbol w\in H^1(\Omega)^2$, 可应用恒等式
\eqref{eq:navier-upwind-integration-by-parts}:
\[
\widetilde a(u_h;v_h-z,\boldsymbol w)
=-a_1(u_h;Pw_h,v_h-z),
\]
所以 Lemma~\ref{lem:navier-upwind-form-bounded} 和
\eqref{eq:navier-upwind-smoothing-regularity} 给出
\[
\begin{aligned}
|\widetilde a(u_h;v_h-z,\boldsymbol w)|
&\leq C_1\|\boldsymbol u_h\|_{0,4,\Omega}|Pw_h|
\|\boldsymbol v_h-\boldsymbol z\|_{0,4,\Omega}\\
&\leq C_1\|\boldsymbol u_h\|_{0,4,\Omega}|w_h|
\|\boldsymbol v_h-\boldsymbol z\|_{0,4,\Omega}.
\end{aligned}
\]
类似地,
\[
\begin{aligned}
|a_1(u_h;v_h-z,\boldsymbol w_h-\boldsymbol w)|
&\leq C_2\|\boldsymbol u_h\|_{0,4,\Omega}
\left(\sum_{\kappa\in\mathcal T_h}|\phi_h-\alpha|_{2,\kappa}^2\right)^{1/2}
\|\boldsymbol w_h-\boldsymbol w\|_{0,4,\Omega}\\
&\leq C_3h^{1/2}\|\boldsymbol u_h\|_{0,4,\Omega}|w_h|
\left(\sum_{\kappa\in\mathcal T_h}|\phi_h-\alpha|_{2,\kappa}^2\right)^{1/2},
\end{aligned}
\]
这里使用了 \eqref{eq:navier-upwind-mesh-dependent-norm} 和
Lemma~\ref{lem:navier-upwind-projection-mesh-norm}. 因此
\[
|a_1(u_h;v_h-z,\boldsymbol w_h-\boldsymbol w)|
+|\widetilde a(u_h;v_h-z,\boldsymbol w)|
\leq C_4h^{1/2}\|\boldsymbol u_h\|_{0,4,\Omega}|w_h|
[\phi_h-\alpha]_4.
\]
最后, 与 Lemma~\ref{lem:navier-upwind-form-bounded} 中相同地导出

\MathBookSourcePage{359}
\[
\begin{aligned}
|a_2(u_h;v_h-z,\boldsymbol w_h-\boldsymbol w)|
&\leq C_5h^{1/2}\|\boldsymbol u_h\|_{0,4,\Omega}
[\phi_h-\alpha]_4[\chi_h-\chi]_2\\
&\leq C_6h^{1/2}\|\boldsymbol u_h\|_{0,4,\Omega}|w_h|
[\phi_h-\alpha]_4.
\end{aligned}
\]
汇集这些不等式即得 \eqref{eq:navier-upwind-consistency-bound}.
\end{Proof}

\setcounter{statement}{0}
\begin{Remark}\label{rem:navier-upwind-gh-lipschitz}
由 Lemma~\ref{lem:navier-upwind-form-bounded} 和
Lemma~\ref{lem:navier-upwind-boundary-almost-lipschitz} 的论证可知,
映射 $G_h$ 在 $V_h$ 上 Lipschitz 连续:
\begin{equation}\label{eq:navier-upwind-gh-lipschitz}
\|G_h(u_h)-G_h(u_h^*)\|_*
\leq C(|u_h|+|u_h^*|)|u_h-u_h^*|
\quad\forall u_h,u_h^*\in V_h,
\end{equation}
其中常数 $C$ 与 $h$ 无关.
\end{Remark}

现在可以定义算子 $\nabla G_h(u_h)$, 即尽可能逼近形式
$\widetilde a(\mathord\cdot;\mathord\cdot,\mathord\cdot)$ 的``导数''.
显然, 最简单的选择是把
$a_2^{z_h}(\mathord\cdot;\mathord\cdot,\mathord\cdot)$ 对 $z_h$ 的依赖线性化.
因此作如下定义.

\setcounter{statement}{0}
\begin{Definition}\label{def:navier-upwind-generalized-gradient}
对所有 $u_h\in V_h$, 算子
$\nabla G_h(u_h)\in\mathcal L(V_h;Y_h)$ 定义为
\[
\langle\nabla G_h(u_h)\cdot v_h,w_h\rangle
=\widetilde a^{u_h}(u_h;v_h,\boldsymbol w_h)
+\widetilde a^{u_h}(v_h;u_h,\boldsymbol w_h)
\quad\forall\boldsymbol w_h=\operatorname{curl}\chi_h,
\quad\chi_h\in\Phi_h.
\]
\end{Definition}

显然, $\nabla G_h(u_h)$ 是从 $V_h$ 到 $Y_h$ 的线性算子, 且
Lemma~\ref{lem:navier-upwind-form-bounded} 蕴含
\[
\|\nabla G_h(u_h)\cdot v_h\|_*
\leq C|u_h||v_h|
\quad\forall u_h,v_h\in V_h,
\]
其中常数 $C$ 与 $h$ 无关.

\setcounter{statement}{1}
\begin{Remark}\label{rem:navier-upwind-gradient-almost-lipschitz}
注意, 关于 $u_h$, $\nabla G_h(u_h)$ ``几乎 Lipschitz 连续''. 事实上,
\[
\begin{aligned}
&|\langle(\nabla G_h(u_h)-\nabla G_h(u_h^*))\cdot v_h,w_h\rangle|\\
={}&|\widetilde a^{u_h}(u_h;v_h,\boldsymbol w_h)
-\widetilde a^{u_h^*}(u_h^*;v_h,\boldsymbol w_h)\\
&\qquad+\widetilde a^{u_h}(v_h;u_h,\boldsymbol w_h)
-\widetilde a^{u_h^*}(v_h;u_h^*,\boldsymbol w_h)|\\
\leq{}&|\widetilde a^{u_h}(u_h-u_h^*;v_h,\boldsymbol w_h)|
+|\widetilde a^{u_h}(u_h^*;v_h,\boldsymbol w_h)
-\widetilde a^{u_h^*}(u_h^*;v_h,\boldsymbol w_h)|\\
&+|\widetilde a^{u_h}(v_h;u_h-u_h^*,\boldsymbol w_h)|
+|\widetilde a^{u_h}(v_h;u_h^*,\boldsymbol w_h)
-\widetilde a^{u_h^*}(v_h;u_h^*,\boldsymbol w_h)|\\
\leq{}&C_1|v_h||w_h|[|u_h-u_h^*|
+h^{1/2}(\|\boldsymbol u_h-\boldsymbol u_h^*\|_{0,4,\Omega}+|u_h^*|)],
\end{aligned}
\]
这里应用了 Lemmas~\ref{lem:navier-upwind-boundary-almost-lipschitz} 和
\ref{lem:navier-upwind-consistency-bound}. 因而
\[
\|\nabla G_h(u_h)-\nabla G_h(u_h^*)\|_{\mathcal L(V_h;Y_h)}
\leq C_2(|u_h-u_h^*|+h^{1/2}|u_h^*|).
\]
当然, 也可导出右端含有 $h^{1/2}|u_h|$ 的类似上界.
\end{Remark}

最后, 为应用 Theorem~\ref{thm:navier-branch-nondifferentiable-approximation},
必须把由 \eqref{eq:navier-upwind-nonlinearity} 定义的算子 $G$ 与
对应于三线性形式 $a_1(\mathord\cdot;\mathord\cdot,\mathord\cdot)$ 的算子联系起来.

\MathBookSourcePage{360}
当 $\boldsymbol u=\operatorname{curl}\psi$ 且
$\psi\in\Phi_s\cap H^2(\Omega)$ 时, 公式
\eqref{eq:navier-centered-stream-convection-curl} 给出
\[
(u_j\partial u_i/\partial x_j,v_i)
=(\operatorname{curl}\boldsymbol u\operatorname{grad}\psi,\boldsymbol v)
\quad\forall\boldsymbol v=\operatorname{curl}\phi,
\quad\phi\in\Phi_s.
\]
因此, 只要 $\psi$ 具有 $H^2$ 正则性,
\[
a_1(u;u,\boldsymbol v)=(\omega\operatorname{grad}\psi,\boldsymbol v).
\]
当然, 当 $u$ 仅属于 $X$ 时, 这个等式不成立, 但可通过正则化函数 $u$ 把 $G$ 延拓到 $X$.
因此对所有 $u\in X$ 定义
\begin{equation}\label{eq:navier-upwind-regularized-nonlinearity}
\langle G(u),\boldsymbol v\rangle
=a_1(u;Pu,\boldsymbol v)-\langle\boldsymbol f,\boldsymbol v\rangle
\quad\forall\boldsymbol v\in L^s(\Omega)^2.
\end{equation}
由于对所有 $u\in V$ 都有 $Pu=u$, 当 $u\in V$ 且
$\boldsymbol v=\operatorname{curl}\phi$, $\phi\in\Phi_s$ 时,
该定义与 \eqref{eq:navier-upwind-nonlinearity} 一致. 此外容易验证,
\eqref{eq:navier-upwind-regularized-nonlinearity} 定义了一个
$C^\infty$ 映射 $G:X\to Y$, 其导数
$DG(u)\in\mathcal L(X;Y)$ 为
\begin{equation}\label{eq:navier-upwind-regularized-derivative}
\langle DG(u)\cdot v,\boldsymbol w\rangle
=a_1(u;Pv,\boldsymbol w)+a_1(v;Pu,\boldsymbol w)
\quad\forall\boldsymbol w\in L^s(\Omega)^2.
\end{equation}
最后可以证明: $u$ 是 \eqref{eq:navier-upwind-continuous-problem} 的非奇异解,
当且仅当 $u$ 也是采用 \eqref{eq:navier-upwind-regularized-nonlinearity}
定义 $G$ 的 Navier--Stokes 问题的非奇异解.

下面检验 Theorem~\ref{thm:navier-branch-nondifferentiable-approximation} 的假设.
由 \eqref{eq:navier-upwind-discrete-stokes} 立即得到
\begin{equation}\label{eq:navier-upwind-th-bound}
\|T_hl\|_X\leq(1+C_s)^2\|l\|_h
\quad\forall l\in Y_h,
\end{equation}
其中 $C_s$ 表示 \eqref{eq:navier-upwind-uniform-w1t} 中的常数. 其次,
Subsection~\ref{subsec:navier-centered-stream-vorticity} 已经检验了
\eqref{eq:navier-branch-th-converges}: 由于正则性条件
\eqref{eq:navier-centered-stream-dual-regularity},
$T_h\boldsymbol f=(\operatorname{curl}\psi_h,\omega_h)$ 满足误差估计
\begin{equation}\label{eq:navier-upwind-stokes-error}
|\psi-\psi_h|_{1,s,\Omega}+\|\omega-\omega_h\|_{0,\Omega}
\leq Ch^\alpha\|\boldsymbol f\|_{0,t,\Omega},
\end{equation}
其中 $r\leq t<2$, $1/\gamma+1/t=1$,
$l=1$ 时 $\alpha=1/\gamma$, 而 $l\geq2$ 时 $\alpha=2/\gamma$.
类似地, 算子 $\pi_h$ 的逼近性质已在
Subsection~\ref{subsec:incompressible-further-error-refinement} 中导出.
事实上, 对 $v=(\operatorname{curl}\phi,\theta)\in V$, 取
\begin{equation}\label{eq:navier-upwind-projection-operator}
\pi_hv=(\operatorname{curl}(\mathring P_h\phi),\theta_h)\in V_h,
\end{equation}
即 $\theta_h$ 由
\[
(\theta_h,\mu_h)
=(\operatorname{curl}(\mathring P_h\phi),\operatorname{curl}\mu_h)
\quad\forall\mu_h\in\Theta_h
\]
确定. 于是
\[
\begin{aligned}
\|\pi_hv-v\|_X\leq{}&
|\mathring P_h\phi-\phi|_{1,s,\Omega}
+2\inf_{\mu_h\in\Theta_h}\|\mu_h-\theta\|_{0,\Omega}\\
&+\sup_{\mu_h\in\Theta_h}
\frac{(\operatorname{curl}(\mathring P_h\phi-\phi),
\operatorname{curl}\mu_h)}{\|\mu_h\|_{0,\Omega}}.
\end{aligned}
\]
因此 Lemma~\ref{lem:incompressible-further-boundary-element-bound},
Theorem~\ref{thm:appendix-ritz-projection-lp-error} 和标准稠密性论证给出
\[
\lim_{h\to0}\|\pi_hv-v\|_X=0
\quad\forall v\in V.
\]
